% !TEX root = ../main.tex

\chapter{Presentation du projet}

\section{Contexte general du projet}

Le domaine financier connait une evolution importante grace a la transformation
numerique, a l'exploitation massive des donnees et aux avancees de l'intelligence
artificielle. Les investisseurs, qu'ils soient particuliers ou professionnels, sont
aujourd'hui confrontes a une quantite tres importante d'informations provenant de
sources variees : rapports financiers, donnees de marche, actualites economiques,
communiques d'entreprises et documents reglementaires.

Malgre cette richesse informationnelle, la prise de decision reste complexe. Les
donnees sont souvent dispersees, difficiles a analyser rapidement et parfois
contradictoires. En particulier, il peut exister un ecart entre le discours public ou
marketing d'une entreprise et les risques reellement mentionnes dans ses documents
officiels. Les rapports 10-K publies aupres de la SEC representent une source fiable
pour identifier ces risques, mais leur exploitation manuelle demande beaucoup de temps
et de competences techniques.

Dans ce contexte, l'utilisation de l'intelligence artificielle permet d'automatiser
l'analyse des documents financiers, de detecter les signaux faibles et d'aider les
investisseurs a prendre des decisions plus argumentees. C'est dans cette perspective
que s'inscrit le projet \textbf{FinPulse}.

\section{Presentation du projet}

FinPulse est une plateforme intelligente d'analyse financiere basee sur les donnees
SEC et l'intelligence artificielle. Elle vise a accompagner l'investisseur dans
l'evaluation de ses idees d'investissement en s'appuyant sur des sources officielles,
notamment les rapports 10-K des entreprises cotees.

La plateforme combine plusieurs modules : un pipeline d'ingestion des donnees
financieres, une base de donnees vectorielle, un assistant conversationnel, un moteur
de strategie, un systeme de surveillance automatisee et une watchlist en temps reel.
L'objectif est de transformer des documents financiers complexes en informations
exploitables, contextualisees et faciles a interpreter.

\subsection{Problematique}

Les investisseurs prennent souvent leurs decisions a partir d'informations publiques,
d'actualites recentes ou d'arguments strategiques personnels. Cependant, ces elements
ne suffisent pas toujours a mesurer la coherence d'une decision d'investissement.
Plusieurs difficultes se posent :

\begin{itemize}
    \item les rapports financiers officiels sont longs, techniques et difficiles a
    analyser manuellement ;
    \item les risques importants peuvent etre caches dans des sections detaillees des
    rapports 10-K ;
    \item le discours marketing d'une entreprise peut ne pas refleter totalement sa
    situation reglementaire ;
    \item les investisseurs ont besoin d'indicateurs clairs pour comparer les preuves,
    les contradictions et les signaux de risque ;
    \item le suivi continu des entreprises et des strategies demande une automatisation
    fiable.
\end{itemize}

La problematique principale du projet peut donc etre formulee ainsi :
\textit{comment concevoir une plateforme intelligente capable d'analyser les donnees
financieres officielles, d'evaluer la coherence d'une strategie d'investissement et
d'aider l'utilisateur a suivre les risques en temps reel ?}

\subsection{Objectifs du projet}

Le projet FinPulse poursuit plusieurs objectifs fonctionnels et techniques :

\begin{itemize}
    \item collecter et exploiter les rapports reglementaires publies par les entreprises
    aupres de la SEC ;
    \item extraire les sections importantes, notamment les risques et les analyses
    financieres ;
    \item utiliser des mecanismes d'intelligence artificielle pour rechercher les
    informations pertinentes dans les documents ;
    \item proposer un assistant conversationnel capable de repondre aux questions de
    l'utilisateur avec des preuves issues des rapports officiels ;
    \item calculer un score de coherence, appele NCI, pour evaluer le niveau de risque
    et la coherence d'une strategie ;
    \item permettre la sauvegarde et la surveillance automatique des strategies
    d'investissement ;
    \item fournir une watchlist en temps reel avec des indicateurs de marche, de
    sentiment et de risque ;
    \item garantir une architecture maintenable, securisee et evolutive.
\end{itemize}

\subsection{Etude de l'existant}

Afin de mieux positionner FinPulse, nous avons etudie plusieurs solutions existantes
utilisees dans le domaine de l'analyse financiere, de la consultation de donnees de
marche et de l'aide a la decision. Le tableau suivant presente quelques exemples
representatifs ainsi que leurs principales limites par rapport aux besoins identifies
dans notre projet.

\begin{table}[htbp]
\centering
\scriptsize
\setlength{\tabcolsep}{3pt}
\renewcommand{\arraystretch}{1.25}
\begin{tabular}{|p{3cm}|p{4cm}|p{6cm}|}
\hline
\textbf{Solution existante} & \textbf{Fonctionnalites principales} & \textbf{Limites observees} \\
\hline
Yahoo Finance &
Consultation des cours boursiers, actualites financieres, graphiques, donnees
fondamentales et informations generales sur les entreprises. &
L'outil reste principalement descriptif. Il fournit des donnees de marche et des
actualites, mais ne permet pas d'analyser automatiquement la coherence entre une
strategie d'investissement et les risques declares dans les rapports 10-K. \\
\hline
Google Finance &
Suivi simple des actions, consultation des prix, evolution historique et informations
de base sur les societes cotees. &
La plateforme offre une vision synthetique mais limitee. Elle ne propose pas
d'analyse approfondie des documents reglementaires, ni de score de risque
personnalise, ni de justification basee sur des sources officielles. \\
\hline
TradingView &
Analyse graphique avancee, indicateurs techniques, alertes de marche, scripts
personnalises et visualisation temps reel. &
La solution est tres orientee analyse technique. Elle ne se concentre pas sur
l'exploitation semantique des rapports SEC et ne detecte pas les contradictions entre
le discours d'une entreprise et ses risques reglementaires. \\
\hline
Bloomberg Terminal &
Acces professionnel a des donnees financieres tres completes, actualites, analyses,
indicateurs de marche et outils de suivi institutionnels. &
La solution est couteuse et destinee principalement aux professionnels. Elle reste
complexe pour un utilisateur non expert et ne fournit pas forcement une analyse RAG
personnalisee a partir d'une idee d'investissement formulee en langage naturel. \\
\hline
\end{tabular}
\caption{Comparaison de solutions existantes orientees marche et analyse}
\label{tab:etude-existant-marche}
\end{table}

\begin{table}[htbp]
\centering
\scriptsize
\setlength{\tabcolsep}{3pt}
\renewcommand{\arraystretch}{1.25}
\begin{tabular}{|p{3cm}|p{4cm}|p{6cm}|}
\hline
\textbf{Solution existante} & \textbf{Fonctionnalites principales} & \textbf{Limites observees} \\
\hline
EDGAR SEC &
Acces officiel aux rapports financiers et documents reglementaires des entreprises
cotees aux Etats-Unis. &
La plateforme est fiable mais peu assistee. Les rapports sont disponibles sous forme
de documents longs et techniques, ce qui rend leur analyse manuelle difficile et
chronophage pour l'investisseur. \\
\hline
Seeking Alpha &
Articles d'analyse, opinions d'experts, actualites financieres, notations et
commentaires sur les entreprises. &
Le contenu depend fortement des analystes et contributeurs. Les analyses ne sont pas
toujours directement reliees aux sections exactes des documents SEC et peuvent
contenir une part d'opinion subjective. \\
\hline
MarketWatch &
Actualites financieres, suivi des marches, informations economiques et donnees
boursieres generales. &
La plateforme est utile pour suivre l'actualite, mais elle ne propose pas de moteur
d'analyse personnalise capable de surveiller une strategie et de declencher des
alertes basees sur l'evolution du score de coherence. \\
\hline
\end{tabular}
\caption{Comparaison de solutions existantes orientees documents et actualites}
\label{tab:etude-existant-documents}
\end{table}

D'apres cette comparaison, les solutions existantes repondent chacune a une partie du
besoin : consultation des cours, analyse technique, acces aux documents officiels ou
lecture d'actualites. Cependant, elles ne combinent pas dans un meme outil l'analyse
semantique des rapports SEC, la recherche augmentee par generation, le calcul d'un
score de coherence, la personnalisation selon la strategie de l'utilisateur et la
surveillance en temps reel. C'est cette complementarite que FinPulse cherche a
apporter.

\subsection{Solution proposee}

La solution proposee est une plateforme nommee \textbf{FinPulse}. Elle repose sur une
architecture moderne composee d'une interface web, d'un backend d'orchestration, d'un
service d'intelligence artificielle et d'une base de donnees centralisee.

FinPulse permet a l'utilisateur de soumettre une idee ou une strategie
d'investissement en langage naturel. Le systeme analyse cette demande, recherche les
informations pertinentes dans les rapports 10-K, calcule un score de coherence et
genere une reponse argumentee. Lorsque l'utilisateur souhaite suivre une entreprise,
la plateforme peut sauvegarder la strategie et declencher des alertes en cas de
variation significative du risque ou d'apparition de nouveaux signaux critiques.

La solution integre egalement une watchlist dynamique permettant de suivre les
entreprises favorites, les variations de marche, les actualites pertinentes et
l'evolution du score NCI.

\section{Methodologie et demarche suivie}

\subsection{Approche methodologique}

La realisation du projet a suivi une demarche progressive, allant de l'analyse du
besoin jusqu'a la conception technique de la solution. Cette approche a permis de
structurer le travail en plusieurs etapes :

\begin{enumerate}
    \item analyse du contexte et identification de la problematique ;
    \item definition des besoins fonctionnels et non fonctionnels ;
    \item etude des solutions existantes et identification de leurs limites ;
    \item conception de l'architecture globale de la plateforme ;
    \item modelisation UML des cas d'utilisation, des classes et des interactions ;
    \item choix des technologies adaptees au projet ;
    \item planification des modules a developper et repartition des responsabilites.
\end{enumerate}

Cette demarche permet de garantir une meilleure maitrise du projet et de reduire les
risques techniques lors de l'implementation.

\subsection{Outils de developpement et de communication}

Plusieurs outils ont ete utilises pour organiser, developper et suivre le projet :

\begin{itemize}
    \item \textbf{Visual Studio Code} et \textbf{IntelliJ IDEA} pour le developpement ;
    \item \textbf{Git} et \textbf{GitHub} pour la gestion du code source et le travail
    collaboratif ;
    \item \textbf{Spring Boot} pour le backend et l'orchestration metier ;
    \item \textbf{Angular} pour le developpement de l'interface utilisateur ;
    \item \textbf{Python} et \textbf{FastAPI} pour le traitement des donnees et les
    services d'intelligence artificielle ;
    \item \textbf{PostgreSQL} et \textbf{PGVector} pour la persistance des donnees et le
    stockage vectoriel ;
    \item \textbf{Redis} pour le cache et l'amelioration des performances ;
    \item \textbf{Docker} pour la containerisation des services ;
    \item des outils de communication et de suivi tels que Google Meet, WhatsApp et
    Trello.
\end{itemize}

\section{Organisation et planification du projet}

\subsection{Equipe du projet}

Le projet a ete realise par une equipe composee de quatre membres. La repartition des
taches a ete faite selon les competences et les responsabilites de chacun, afin de
couvrir les differents aspects du projet : ingestion des donnees, intelligence
artificielle, backend, frontend, base de donnees, conception UML et redaction du
rapport.

Cette organisation a permis de travailler en parallele sur plusieurs modules tout en
assurant une coordination reguliere entre les membres de l'equipe. Les reunions de
suivi ont permis de discuter de l'avancement, de resoudre les difficultes techniques
et d'ajuster les priorites.

\subsection{Deroulement du projet}

Le deroulement du projet peut etre divise en plusieurs phases principales :

\begin{itemize}
    \item \textbf{Phase d'analyse} : comprehension du sujet, definition du besoin et
    formulation de la problematique ;
    \item \textbf{Phase de conception} : elaboration des diagrammes UML, definition de
    l'architecture et choix des technologies ;
    \item \textbf{Phase de developpement} : implementation des differents modules de la
    plateforme ;
    \item \textbf{Phase d'integration} : connexion entre le frontend, le backend, les
    services d'IA et la base de donnees ;
    \item \textbf{Phase de test} : verification du bon fonctionnement des
    fonctionnalites principales ;
    \item \textbf{Phase de documentation} : redaction du rapport et preparation de la
    presentation finale.
\end{itemize}

\section{Impact attendu}

FinPulse vise a apporter une valeur ajoutee importante dans le domaine de l'aide a la
decision financiere. La plateforme permet de reduire le temps necessaire a l'analyse
des documents reglementaires, d'ameliorer la comprehension des risques et de fournir
des recommandations plus transparentes grace a la citation des sources officielles.

L'impact attendu concerne aussi bien l'utilisateur final que l'aspect technique. Pour
l'investisseur, FinPulse offre un outil d'analyse plus clair, plus rapide et plus
personnalise. Sur le plan technique, le projet demontre l'interet de combiner les
architectures modernes, les bases vectorielles et les modeles d'intelligence
artificielle pour exploiter des donnees financieres complexes.

\section{Conclusion}

Ce chapitre a presente le cadre general du projet FinPulse, sa problematique, ses
objectifs, l'etude de l'existant ainsi que la solution proposee. Il a egalement decrit
la methodologie adoptee, les outils utilises, l'organisation de l'equipe et l'impact
attendu de la plateforme.

Cette presentation constitue une base pour la suite du rapport, qui detaille la
conception fonctionnelle et technique de FinPulse, ainsi que les choix architecturaux
retenus pour repondre aux besoins identifies.
